\documentclass[11pt]{article}
\usepackage{amsmath,amssymb}
\usepackage[margin=1in]{geometry}
\usepackage{booktabs}
\usepackage{hyperref}
\usepackage{xcolor}
\hypersetup{colorlinks=true, linkcolor=blue!50!black, urlcolor=blue!50!black}

\newcommand{\sat}[2]{[#1]\!:\!#2^{\infty}}
\newcommand{\Sbar}{\bar S}
\newcommand{\code}[1]{\texttt{#1}}

\title{Rosenfeld--Gr\"obner regularization, base-field saturation,\\
       and the bad locus: why the hydrogen ansatz\\ acquires a spurious
       component $(a_0,a_1)$}
\author{}
\date{\today}

\begin{document}
\maketitle

\noindent
A companion note to \emph{A New Method of Solving PDEs}, documenting a subtlety
that surfaced in the computational backing. The projection script
\code{joca.sage} decomposes the hydrogen solution locus into \textbf{five}
prime strata. The Rosenfeld--Gr\"obner--augmented variant \code{joca-rg.sage},
\emph{as originally written}, returned only \textbf{four}, replacing two genuine
strata by a single spurious component $(a_0,a_1)$. This note explains exactly
why, reconciles it with the core theorem of the Rosenfeld--Gr\"obner algorithm,
exhibits the saturation explicitly via \code{normal\_form}, and records which
script does what.

\section{Setup: the ansatz and the projection}

We seek solutions of the zero-energy hydrogen Schr\"odinger equation
$-\tfrac12\nabla^2\Psi-\tfrac1r\Psi=E\Psi$ (cleared of $1/r$) of the form
$\Psi=\Psi(v)$, where $v=v_1x+v_2y+v_3z+v_4r$, $r=\sqrt{x^2+y^2+z^2}$, and $\Psi$
satisfies a second-order linear ODE with affine coefficients. As a differential
system this is the \emph{ansatz}
\begin{equation}
\label{eq:ansatz}
\begin{aligned}
  &\Psi_x-\Psi' v_x,\;\; \Psi_y-\Psi' v_y,\;\; \Psi_z-\Psi' v_z,\quad
   \Psi'_x-\Psi'' v_x,\;\; \Psi'_y-\Psi'' v_y,\;\; \Psi'_z-\Psi'' v_z,\\
  &\boxed{(a_0+a_1v)\,\Psi'' + (b_0+b_1v)\,\Psi' + (c_0+c_1v)\,\Psi},\quad
   v-(v_1x+v_2y+v_3z+v_4r),\quad r^2-x^2-y^2-z^2,
\end{aligned}
\end{equation}
with the eleven constants $E,v_1,\dots,v_4,a_0,a_1,b_0,b_1,c_0,c_1$ as
parameters. The \textbf{projection method} (\code{joca.sage}) reduces the PDE
modulo \eqref{eq:ansatz} (Ritt reduction), collects the coefficients of the
remainder as polynomials in the constants, and computes the minimal associated
primes of the resulting ideal. The redundancy locus---the set of constants on
which the PDE is implied by the ansatz---is

\begin{equation}
\label{eq:five}
\begin{aligned}
P_1 &= (a_0,v_4,v_3,v_2,v_1), &
P_2 &= (c_1,c_0,b_1,b_0,a_1,a_0), &
P_3 &= (v_1^2{+}v_2^2{+}v_3^2,\,a_1,a_0,v_4),\\[2pt]
P_4 &= (v_4c_0{-}b_0,\,Ec_0{-}v_4c_1,\,Eb_0{-}v_4^2c_1,\,b_1,\,a_1{-}\tfrac12 b_0,\,a_0,v_3,v_2,v_1),\\[2pt]
P_5 &= (v_4c_0{-}b_0,\,v_1^2{+}v_2^2{+}v_3^2{-}v_4^2,\,c_1,b_1,\,a_1{-}b_0,\,a_0,\,E).
\end{aligned}
\end{equation}
These five primes are verified (§1.12 cross-check) to be exactly the cells on
which the PDE reduces to $0$ modulo the ansatz.

\section{The RG-augmented variant and the spurious \texorpdfstring{$(a_0,a_1)$}{(a0,a1)}}

The general algorithm's hypothesis (3) asks that the ansatz be a coherent,
squarefree \emph{regular differential system}. \code{joca-rg.sage} discharges it
by running \code{RosenfeldGroebner} over the coefficient field
$\mathbb{Q}(E,v_1,\dots,c_1)$ (the constants must be moved into the base field
for RG to terminate---see \code{rg\_basefield.py} and the cliff in
\code{joca-rg-orderly.sage}). RG returns a single regular component. \emph{Its
equations are not the input.} In particular the autoreduction
\emph{squares the ODE's initial}: the regularized ODE is
\begin{equation}
\label{eq:squared}
   (a_0+a_1v)\cdot\bigl[(a_0+a_1v)\Psi''+(b_0+b_1v)\Psi'+(c_0+c_1v)\Psi\bigr].
\end{equation}
If one then reduces the PDE against these regularized \emph{equations} and
prime-decomposes, the decomposition collapses to four primes:
$P_1,\;(a_1,a_0),\;P_4,\;P_5$, with the single coarse $(a_1,a_0)$ standing in for
$P_2$ and $P_3$.

\paragraph{$(a_0,a_1)$ is spurious, not a coarsening.} It is not that $P_2,P_3$
``merge''; the cell $(a_0,a_1)$ strictly \emph{over}-claims. Indeed
$V(a_0,a_1)\supsetneq V(P_2)\cup V(P_3)$, and the true redundancy ideal
$I=\bigcap_i P_i$ contains the generator $v_4 b_1$, which does \emph{not} vanish
at a generic point of $a_0=a_1=0$ (take $v_4,b_1\neq0$). So the PDE is genuinely
\emph{not} redundant on generic $a_0=a_1=0$: the reported $(a_0,a_1)$ is an
artifact of the reduction, not a solution stratum.

\section{Diagnosis: the missing saturation}

A regular differential chain $C$ does not represent the bare ideal $[C]$; it
represents the \emph{saturated} ideal $\sat{C}{H_C}$, where
$H_C$ is the product of the initials and separants of $C$. This is the content
of Theorem~28 of Boulier--Lazard--Ollivier--Petitot, \emph{Computing
representations for radicals of finitely generated differential ideals}: RG
computes regular systems $(A_i,S_i)$ with
\begin{equation}
\label{eq:thm28}
   \sqrt{\sat{P_0}{S_0}} \;=\; \sat{A_1}{S_1}\,\cap\,\cdots\,\cap\,\sat{A_n}{S_n}.
\end{equation}
The final autoreduction rewrites the equations \emph{and} grows the inequation
set: at each step $\bar p = p\,\text{full-rem}\,\varphi p'$ and
$S_{k+1}=S_k\cup\{i_{\bar p},s_{\bar p}\}$ (invariant I4: $i_p,s_p\in S$ for every
$p$). So the injected initial $(a_0+a_1v)$ in \eqref{eq:squared} is exactly
\emph{compensated} by $\Sbar$ containing it: the saturation $:\Sbar^\infty$
divides it back out, and the represented ideal $\sat{\bar A}{\Sbar}$ is unchanged.

\medskip
\noindent\textbf{This pinpoints the bug.} Boulier's object is the pair
$(\bar A,\Sbar)$ representing $\sat{\bar A}{\Sbar}$. The method
\code{equations()} returns only $\bar A$. Reducing the PDE by $\bar A$ and
prime-decomposing computes membership in $[\bar A]$---\emph{dropping the
$:\Sbar^\infty$}. The spurious $(a_0,a_1)$ is nothing but the locus
$\Sbar=0$ (i.e.\ $a_0+a_1v=0\iff a_0=a_1=0$) that the saturation was meant to
delete. In one line:
\[
  \underbrace{\sat{\bar A}{\Sbar}}_{\text{radical, 5 primes}}
  \quad\text{vs.}\quad
  \underbrace{[\bar A]}_{\text{radical }+\ \{\Sbar=0\}=(a_0,a_1)}.
\]

\section{The saturation made explicit: \texorpdfstring{\code{normal\_form}}{normal\_form}}

The chain's own \emph{saturation-aware} reduction, \code{normal\_form}, displays
the $:\Sbar^\infty$ directly. It returns the normal form of the PDE as a
\textbf{fraction} $\mathrm{NUM}/\mathrm{DEN}$, where the denominator is the
chain separant
\begin{equation}
\label{eq:den}
  \mathrm{DEN}\;=\;\Sbar\;=\;\bigl(a_0+a_1v_\parallel\bigr)^2-a_1^2v_4^2\,(x^2{+}y^2{+}z^2),
  \qquad v_\parallel:=v_1x+v_2y+v_3z,
\end{equation}
the reduced $(a_0+a_1v)^2$ modulo $r^2=x^2+y^2+z^2$---equivalently the norm
$(a_0+a_1v)(a_0+a_1\bar v)$ over the quadratic relation, where
$\bar v=v_\parallel-v_4r$ reflects the radial part. The membership condition
is $\mathrm{NUM}=0$ \emph{with} $\mathrm{DEN}\neq0$---one does not zero the
denominator, since $\Sbar$ is declared a non-zero-divisor. Crucially,
\begin{center}
\fbox{$\mathrm{DEN}$ vanishes exactly on $a_0=a_1=0$,}
\end{center}
so $a_0=a_1=0$ is a \textbf{pole} of the saturated reduction. There the normal
form is $0/0$: undefined. \code{normal\_form} therefore says \emph{nothing}
about the bad locus---neither redundant nor not. The locus
$\{\mathrm{NUM}=0,\ \mathrm{DEN}\neq0\}$ is the redundancy \emph{off}
$a_0=a_1=0$, namely the three \emph{generic} strata $P_1,P_4,P_5$. The two
strata $P_2,P_3$ live inside the pole and are invisible to it.

\section{Why the raw projection recovers all five}

The difference between the two routes is what the ODE degenerates \emph{to} on
the bad locus $a_0=a_1=0$:
\begin{itemize}
  \item the \textbf{regularized} ODE \eqref{eq:squared} becomes $0=0$---useless
        as a reduction rule, hence the genuine singularity of the chain there;
  \item the \textbf{raw} ODE $(a_0+a_1v)\Psi''+(b_0+b_1v)\Psi'+(c_0+c_1v)\Psi$
        drops to the \emph{first-order} rule $(b_0+b_1v)\Psi'+(c_0+c_1v)\Psi$,
        which is still usable: the leader slides from $\Psi''$ down to $\Psi'$,
        and reduction continues.
\end{itemize}
The raw ansatz thus covers the whole parameter space uniformly, and its
plain Ritt reduction (\code{joca.sage}) detects $P_2,P_3$ along with the rest.
This is why the corrected \code{joca-rg.sage} uses RG \emph{only} as the
coherence check (discharging hypothesis (3)) and then projects against the
\emph{raw} ansatz, recovering all five primes \eqref{eq:five}.

\paragraph{Reconciliation with Theorem~28.} There is no contradiction. The ideal
$\sat{\bar A}{\Sbar}$ equals the radical (so RG is a no-op \emph{on the ideal}).
But the \emph{saturated reduction} of that ideal carries $\Sbar$ in the
denominator, so its parametric projection is blind to $\Sbar=0$: it returns the
three generic strata. The \emph{raw} representation of the same radical has no
such denominator, so its projection covers $\Sbar=0$ too---all five. To recover
the bad-locus strata through the saturated route one must recurse into the pole,
running \code{RosenfeldGroebner} over the base field \emph{modulo} $(a_0,a_1)$
and taking the union (the comprehensive-triangular-decomposition scheme of
Chen--Moreno-Maza and Montes). The raw projection collapses that recursion into
a single pass.

\section{The scripts and their outputs}

All scripts live in \code{\textasciitilde/Papers/NewMethod/} and run under
SageMath~10.7 with Fran\c{c}ois Boulier's \code{DifferentialAlgebra} 5.3
(\code{sage -python} for the \code{.py} files).

\begin{center}
\renewcommand{\arraystretch}{1.35}
\begin{tabular}{@{}p{0.27\linewidth}p{0.13\linewidth}p{0.52\linewidth}@{}}
\toprule
\textbf{Script} & \textbf{Output} & \textbf{Role} \\
\midrule
\code{joca.sage} & 5 primes, $\approx$5\,s &
  The projection method: reduce the PDE modulo the raw ansatz, collect
  coefficients, prime-decompose. The faithful decomposition \eqref{eq:five}; the
  one the paper prints. \\
\code{joca-rg.sage} & 5 primes, $\approx$26\,s &
  Adds the hypothesis-(3) coherence check: \code{RosenfeldGroebner} over the
  base field returns one regular component, then the script projects against the
  \emph{raw} ansatz. Agrees with \code{joca.sage}. (Its header documents why one
  must \emph{not} reduce against the regularized chain.) \\
\code{joca-normalform.sage} & fraction \mbox{+} 4/3 primes &
  This note's companion. Calls the chain's \code{normal\_form}; prints
  $\mathrm{NUM}/\mathrm{DEN}$ with $\mathrm{DEN}=\Sbar$; shows $\mathrm{DEN}=0$ on
  $a_0=a_1=0$; decomposes $\{\mathrm{NUM}=0\}$ (4 primes, the $(a_0,a_1)$ being
  the pole) and, after saturating by $(a_0,a_1)$, the 3 generic strata. \\
\code{rg\_basefield.py} & 1 component, $\approx$16\,s &
  Minimal standalone demo (no Sage): builds the ring, moves the constants into
  the base field, runs \code{RosenfeldGroebner(ansatz, basefield=F)}, shows it
  terminates with one regular component. \\
\code{joca-rg-orderly.sage} & times out &
  Cliff artifact: RG with the constants left as \emph{ring variables} under an
  orderly ranking. Demonstrates the blow-up is ranking-independent and that the
  base field is what makes RG terminate. \\
\bottomrule
\end{tabular}
\end{center}

\noindent
A related experiment lives in the separate repository
\code{\textasciitilde/parametric-rg} (a from-scratch port of Fakouri--Rahmany--
Basiri's \emph{parametric} Rosenfeld--Gr\"obner): \code{examples/joca-prg.sage}
runs the §1.12 cross-check (all five primes are PDE-redundancy cells) and a
direct parametric-RG attempt on the ansatz that does not terminate---the same
radius-driven coefficient swell, there in the differential reduction rather than
in a saturation. See that repository's \code{README} for the localized
diagnosis.

\section*{Summary}

The component $(a_0,a_1)$ was never in the radical. It is the bad locus $B$ where
an initial vanishes, re-admitted because the script reduced by $[\bar A]$ instead
of $\sat{\bar A}{\Sbar}$. \code{normal\_form} exhibits the missing $:\Sbar^\infty$
as a literal denominator, with $B$ as its pole; and the genuine sub-strata of $B$
are recoverable only by the raw projection (whose ODE degenerates to a usable
first-order rule on $B$) or by recursing into $B$. The corrected
\code{joca-rg.sage} therefore uses RG to certify coherence and projects raw,
matching \code{joca.sage}'s five primes.

\begin{thebibliography}{9}
\bibitem{BLOP}
F. Boulier, D. Lazard, F. Ollivier, M. Petitot, Computing representations for
radicals of finitely generated differential ideals, \emph{Appl.\ Algebra Eng.\
Commun.\ Comput.} \textbf{20} (2009) 73--121 (Theorem~28).
\bibitem{BLS}
F. Boulier, F. Lemaire, A. Sedoglavic, On the regularity property of
differential polynomials modulo regular differential chains, \emph{CASC 2011},
LNCS \textbf{6885}, 61--72 (the \code{normal\_form} algorithm).
\bibitem{CTD}
C. Chen \emph{et al.}, Comprehensive triangular decomposition, \emph{CASC 2007},
LNCS \textbf{4770}, 73--101 (the recurse-into-the-bad-locus scheme).
\end{thebibliography}

\vfill
\noindent\rule{0.4\linewidth}{0.4pt}\\
{\small Prepared by Claude (Opus 4.8) for Brent Baccala. /Opus 4.8}

\end{document}
